\section{Product Overview}

\subsection{Product Perspective}

The Medication Safety Assistant is a standalone,
responsive Web prototype that helps patients identify
potential medication risks and manage confirmed
medication schedules. It is not a diagnostic or
prescribing system and does not connect to real
healthcare systems.

The application has two main user roles: Patient
and Healthcare Professional (HCP). Patients manage
their medications, while HCPs use separate accounts
to review safety cases. HCP review is simulated
and does not represent real clinical approval.
Family caregivers may assist patients with permission.

The system consists of five main modules:

\begin{enumerate}
  
    \item \textbf{Medication Intake and Intent Recognition:}
    Identifies user intent through UI selection or
    NLP, and collects medication information
    through OCR, NLP Chatbot, Voice/STT, and
    Manual Input. Users review and confirm
    medication details before further processing.

    \item \textbf{Medication Identification:}
    Identifies Vietnamese medications and active
    ingredients using permitted data sources.
    Uncertain matches remain unresolved.

    \item \textbf{Medication Safety Assessment:}
    Checks supported drug-drug interactions using
    DDInter 2.0 and detects duplicate active
    ingredients. It reports potential risks,
    warnings, and data limitations without
    claiming that a prescription is fully safe.

    \item \textbf{HCP Review:}
    Allows HCPs to review high-risk or uncertain
    cases, request corrections, and record decisions
    through a simulated review workflow.

    \item \textbf{Medication Reminders and Adherence:}
    Supports confirmed medication schedules,
    applies a safety gate before activating new
    reminders, and tracks medication adherence.

\end{enumerate}

The workflow starts in Module 1. For OCR and
Manual Input, users select their intended action
through the interface. For typed and voice requests,
NLP identifies intent from text or speech transcripts.
Unclear requests require clarification.

After user confirmation, Module 2 identifies
medications, and Module 3 assesses supported
safety risks. The system then displays results,
warnings, and data limitations.

Cases requiring HCP review, including severe
interactions, are handled by Module 4.
Missing or unclear medication information
must not be treated as safe.

After safety assessment and any required
HCP review, the system may suggest setting
medication reminders. Users can accept,
decline, or request reminders directly.

Users may confirm and adjust reminder times
as long as they follow a verified prescription
or recorded HCP instructions. Changes that
conflict with these instructions, or cannot
be verified, require HCP review.

Before activating new reminders, Module 5
applies a safety gate. New reminders for
medicines with severe interactions remain
blocked until HCP review permits activation.

Other eligible reminders may be activated.
The system does not automatically stop
existing reminders or prescribed medicines.

\subsection{Assumptions and Dependencies}

\begin{longtable}{|L{0.16\textwidth}|L{0.43\textwidth}|L{0.31\textwidth}|}
\hline
\textbf{Type} &
\textbf{Assumption / Dependency} &
\textbf{Effect if Invalid or Unavailable}\\ \hline

Assumption
& Users can review, correct, and confirm
medication information before processing.
& Unconfirmed information cannot be used
for safety assessment or reminders.
\\ \hline

Dependency
& Vietnamese medication data and reliable
active-ingredient mappings are available
for medication identification.
& Unknown medicines or unclear mappings
remain unresolved and cannot be
assessed automatically.
\\ \hline

Dependency
& DDInter 2.0 provides interaction data
for the identified active ingredients.
& Unsupported interactions cannot be
assessed. Missing results are not
treated as proof of safety.
\\ \hline

Dependency
& OCR can extract medication information
from clear printed or digital prescriptions.
& Users must correct extraction errors
or switch to Manual Input.
\\ \hline

Dependency
& Speech-to-text can convert voice input
into editable text.
& Users can correct the transcript
or switch to typed or manual input.
\\ \hline

Dependency
& NLP can identify user intents and
medication information from text.
& Users must clarify unclear requests
or select a function manually.
\\ \hline

Assumption
& HCP accounts are available for simulated
review of high-risk or uncertain cases.
& Cases awaiting required review remain
pending. Affected new reminders cannot
be activated without the required decision.
\\ \hline

Assumption
& Prescription or HCP instructions provide
enough information to check reminder times.
& Unclear or conflicting schedule changes
require HCP review before activation.
\\ \hline

Dependency
& Web browsers support camera, microphone,
and reminder notifications when needed.
& Users can use Manual Input if needed.
Notifications may be limited or unavailable.
\\ \hline

\end{longtable}